\chapter{Implementazione dei test automatici}
\label{cap:test-implementazione}

\intro{Il presente capitolo descrive l'introduzione di una strategia di testing
automatico nel framework \emph{Sia.Server.Framework}, fino a quel momento privo di
qualsiasi struttura di test versionata. Vengono illustrate le quattro tipologie di
test adottate, descrivendo per ciascuna cosa esercita concretamente, le scelte
tecnologiche, i vantaggi ottenuti e i difetti latenti emersi grazie ai test stessi.}

\section{Contesto}

Prima dell'intervento, l'unica forma di verifica disponibile era la prova manuale
dell'interfaccia. Le conseguenze pratiche erano la mancanza di una rete di sicurezza
in fase di \emph{refactor}, l'impossibilità di modificare i moduli \emph{core}
(\texttt{Sia.Base}, \texttt{Sia.Grid}, \texttt{Sia.Devextreme.Crud}) con fiducia, e
l'assenza di una \emph{baseline} oggettiva per misurare il comportamento atteso al
variare dei dati di configurazione.

L'obiettivo dell'intervento è stato duplice: impostare una convenzione di testing
replicabile su tutti i moduli del framework, e fornire una prima copertura concreta
sui moduli con maggiore superficie pubblica, comprensiva di tutte e quattro le
tipologie di test descritte nelle sezioni seguenti.

\section{Modello a quattro livelli}

La strategia adottata distingue quattro tipologie di test con responsabilità
complementari e nessuna sovrapposizione, organizzate lungo un asse che va dall'unità
più piccola di codice all'esperienza completa dell'utente.

\subsection{Unit test backend}

Gli \emph{unit test} backend esercitano il singolo controller con tutte le sue
dipendenze sostituite da \emph{mock}: il servizio di dominio non viene mai invocato
realmente, le chiamate a \emph{SignalR}, al \emph{logger}, al servizio di sessione e
al servizio di permessi sono simulate. Ciò che viene verificato è il comportamento
osservabile del controller a fronte di un \emph{service} configurato per rispondere
in un certo modo: il tipo di \texttt{IActionResult} prodotto, gli argomenti passati
al \emph{service}, l'invocazione di \texttt{IPermissionCheckService.CheckPermission}
con il \texttt{PermissionType} corretto, l'invio di notifiche \emph{SignalR} sulle
mutazioni, le chiamate di log \emph{start}/\emph{ok}/\emph{error}, e simmetricamente
la \emph{non}-invocazione del service e di \emph{SignalR} sui rami di errore.

Questo livello è veloce (l'intera suite gira in pochi secondi), deterministico, e
non richiede né database né rete. Per ogni \emph{endpoint} vengono scritti almeno
tre test: uno sul ramo positivo (verifica risposta, parametri e \emph{side-effect}),
uno sul ramo di validazione (verifica che in presenza di input non valido il service
non venga mai invocato e che il \emph{logger} riceva \texttt{LogBadRequest}), e uno
sul ramo eccezionale (verifica che il controller propaghi correttamente le eccezioni
del service avvolgendole nella forma attesa).

\subsection{Integration test backend}

Gli \emph{integration test} verificano che l'intera \emph{pipeline} funzioni con il
database reale: l'host ASP.NET Core viene avviato in-process tramite
\texttt{WebApplicationFactory}, le richieste HTTP vengono eseguite contro un
\texttt{TestServer} senza aprire \emph{socket} reali, e il flusso attraversa tutti
gli strati (routing, model binding, middleware, controller, service, repository)
fino alle \emph{stored procedure} su SQL Server. Restano \emph{mockati} solo gli
strati di autenticazione e permessi, sostituiti da un \texttt{TestAuthHandler} che
impersona un utente noto e da una implementazione \texttt{AllowAll} di
\texttt{IPermissionCheckService}, per evitare \emph{coupling} con utenti specifici
del database.

Questo approccio esercita ciò che gli \emph{unit test} per costruzione non possono
toccare: il \emph{routing}, il \emph{model binding}, i middleware di autenticazione,
l'esecuzione reale delle query SQL e la correttezza della configurazione
\emph{per-componente} memorizzata nelle tabelle del database
(\texttt{component.Components}, \texttt{component.Layouts}, ecc.).

Sono stati prodotti tre livelli di \emph{integration test}.

\paragraph{Smoke.} Itera su tutti i componenti presenti nel database e per ciascuno 
effettua una chiamata reale all'\emph{endpoint} di lettura
(\texttt{POST /api/grid-core/page} oppure \texttt{GET /api/CrudData/page}),
asserendo che la risposta sia \texttt{200 OK} con \emph{payload} JSON valido. Lo scopo è
esercitare ogni singola configurazione \emph{per-componente} presente in produzione,
intercettando \emph{stored procedure} malformate, \emph{connection name} non
risolvibili o \emph{misconfiguration} di colonne. I componenti legittimamente non
testabili (form action-only, \emph{connection} non disponibili localmente, SP
malformate sul DB di sviluppo) sono dichiarati in una \emph{skip list}, che trasforma esiti rumorosi
in \emph{skipped} tracciati e revisionabili in un unico punto.

\paragraph{Lifecycle.} Verifica il ciclo completo di scrittura su un \emph{set}
curato di componenti rappresentativi: crea un record (\texttt{INSERT}), lo legge
(\texttt{GET}), lo modifica (\texttt{UPDATE}), verifica la modifica, lo elimina
(\texttt{DELETE}) e verifica l'avvenuta eliminazione. Ogni componente è descritto
da un file JSON \emph{embedded} che dichiara \emph{payload} di create, \emph{shape}
attesa, delta di update, modalità di delete e --- per le entità con cancellazione
logica --- i campi attesi dopo il \emph{soft-delete} (es.\ \texttt{active=false}). Il
\emph{framework} di test supporta trasparentemente sia chiavi primarie singole sia
composite (es.\ \texttt{crud-signalr} con PK \texttt{(SenderId, MessageTypeId,
RecepientId)}), sia cancellazione fisica sia logica. L'esecuzione avviene in
\emph{collection} serializzata (\texttt{DisableParallelization=true}) per evitare
\emph{race condition} sulle scritture, e in caso di fallimento intermedio un blocco
\texttt{try/finally} esegue una \texttt{DELETE} \emph{best-effort} sul record creato,
evitando di lasciare residui. Il \emph{set} consolidato copre $12$ pattern
strutturali distinti per il CRUD (tabelle dirette, SP standard, SP con \emph{soft-delete},
PK \texttt{varchar}, PK \texttt{IDENTITY}, PK composite) e $6$ griglie pilota per
\texttt{Sia.Grid}.

\paragraph{Summary.} Specifico di \texttt{Sia.Grid}, esercita l'\emph{endpoint}
\texttt{POST /api/grid-core/summary} responsabile del calcolo aggregato
(\texttt{sum}, \texttt{avg}, \texttt{count}, \texttt{min}, \texttt{max}). La
\emph{theory} è completamente \emph{data-driven}: scansiona a \emph{runtime} le
\emph{perspective} JSON in \texttt{Configurations/perspectives/core/grid/} e
individua le griglie che dichiarano almeno una colonna con
\texttt{summaryColEnable: true} e \texttt{summaryType} valorizzato; per ciascuna
costruisce dinamicamente il \emph{payload} e asserisce \texttt{200 OK} con risposta
serializzabile. L'esito tipico è $9/10$ \emph{green} con $1$ \emph{skipped} dovuto a
una \emph{misconfiguration} (\texttt{summaryType=sum} su colonna \texttt{varchar})
correttamente segnalata e tracciata.

Un episodio significativo emerso durante i \emph{lifecycle test} merita menzione:
tutti i test di scrittura fallivano allo step \texttt{INSERT} con un errore SQL
identico, riferito a \emph{synonym} \texttt{*\_Log} che puntavano a un database di
audit (\texttt{CoreDBLog}) inesistente sul server di sviluppo. L'audit, di fatto,
non era mai stato operativo su quell'ambiente: i \emph{trigger} \texttt{TR\_LOG\_*}
fallivano in fase di risoluzione metadati prima ancora di valutare la propria
guardia \texttt{OBJECT\_ID}. I test hanno reso visibile una situazione preesistente
e latente; la soluzione adottata è stata rimuovere i $316$ \emph{synonym} obsoleti,
rendendo le scritture nuovamente funzionanti.

\subsection{Unit test frontend}

Sul lato \emph{frontend}, la libreria Angular \texttt{sia-grid} è stata scelta come
modulo pilota in simmetria con il corrispettivo backend. Gli \emph{unit test} sono
scritti con \textbf{Vitest}, integrato nativamente con il toolchain Vite già in
uso nel workspace, e verificano in isolamento la logica TypeScript pura senza
\emph{render} reale del template: l'\texttt{HttpClient} è sostituito da \emph{mock},
i \emph{signal input} sono settati direttamente sull'istanza.

Concretamente i test esercitano: la corretta costruzione delle \texttt{HttpRequest}
negli \emph{API client} (\emph{path}, \emph{query parameters}, serializzazione del
\emph{payload}, propagazione degli errori); la logica di stato e di trasformazione
dei servizi sui dati delle colonne, con la gestione della cache di configurazione e
il comportamento sugli \emph{observable}; il comportamento delle \emph{directive}
(ad esempio il calcolo del delta durante un \emph{drag} di resize); i metodi puri
del componente principale \texttt{SiaMatGridComponent} ($\sim$$4900$ righe), come
gli \emph{helper} \texttt{isGroup}, \texttt{hasSelected}, \texttt{getDescription}
con \emph{i18n}, i \emph{signal computed} \texttt{selectionMode},
\texttt{compactViewActive}, \texttt{isRadar}, e il \emph{debounce} di $400$\,ms su
\texttt{applyGlobalFilter} verificato con \emph{fake timer}; la logica delle pagine
di configuratore (\emph{setter} di \texttt{selectedPerspectiveId}, \emph{routing} di
\texttt{manageActionFired}, \emph{insert}/\emph{update} di \texttt{savePerspective}
con flag che impedisce chiamate concorrenti).

Questi test non richiedono né \emph{browser} né \emph{backend} attivo, girano in
pochi secondi e sono deterministici. Resta escluso da questo livello tutto ciò che
richiede il \emph{render} reale del \emph{template} DevExtreme/Material o
l'integrazione con il \emph{backend}: \emph{rendering} delle righe, \emph{sort}
effettivo, \emph{export} Excel, \emph{drag-and-drop} delle colonne, \emph{group} e
\emph{master-detail}, \emph{branch} di \texttt{manageSiaSignalrMessage} con
\emph{match} su \texttt{gridId} --- territorio della suite \emph{end-to-end}.

\subsection{Test end-to-end con Playwright}

I test \emph{end-to-end} (E2E) rappresentano il livello più alto della piramide: il
\emph{frontend} Angular (\texttt{ng serve}) e il \emph{backend} .NET (\texttt{Sia.Devs}) sono
entrambi avviati realmente, il \emph{database} è quello vero, e \textbf{Playwright}
guida un'istanza di Chromium \emph{headless} come farebbe un utente umano. Niente
\emph{mock}: ogni \emph{assert} attraversa l'intera catena reale, dal click sul
menu fino alla \emph{stored procedure} SQL e ritorno.

Il vantaggio è la fedeltà assoluta alla realtà: un singolo \emph{assert} (``la prima
riga della tabella è visibile e contiene del testo'') copre simultaneamente il
routing Angular, la costruzione del \emph{payload} HTTP, l'autenticazione, il
middleware del backend, l'esecuzione della \emph{stored procedure}, la
serializzazione del \emph{response} e il rendering del DOM. Il costo è la velocità
(ogni test richiede secondi, non millisecondi) e la dipendenza dallo stato del
sistema. Per questo i test E2E si concentrano sui flussi utente strategici,
lasciando la copertura della logica di dettaglio agli \emph{unit test}.

La suite pilota copre il componente \texttt{sia-mat-grid} con nove \emph{spec}
distinti, ciascuno focalizzato su un singolo aspetto del comportamento utente:

\begin{description}
  \item[\emph{Smoke}] Apre il menu \emph{Elenco Componenti}, attende che la tabella
    \texttt{table[id="siaTablecomponents"]} sia visibile, verifica la presenza
    della \emph{header row}, conta le righe del \emph{body} e asserisce che siano
    $> 0$, infine controlla che il paginatore Material mostri una \emph{label} di
    \emph{range} valida (es.\ ``$1 - 25$ di $342$''). Funge da \emph{canary}: se
    questo test fallisce, qualsiasi altro test sarebbe inutile.
  \item[\emph{Sorting}] Legge la sequenza dei valori della prima colonna con
    \emph{data field}, clicca sull'\emph{header} per innescare il riordino, attende
    via \texttt{expect.poll} che il \emph{re-render} sia stabile, e verifica che i
    valori siano effettivamente ordinati confrontandoli con \texttt{Array.sort}.
    Poiché \texttt{sia-mat-grid} ha una propria implementazione del sort (non
    \texttt{mat-sort-header}), non ci si può affidare agli attributi ARIA standard:
    la verifica è comportamentale.
  \item[\emph{Filtering}] Imposta un filtro ``inizia con V'' sulla prima colonna
    filtrabile, verifica che tutti i valori visibili rispettino il criterio, poi
    rimuove il filtro e verifica (\texttt{expect.poll}) che il numero di righe torni
    esattamente a quello iniziale. È un \emph{round-trip} completo che esercita sia
    l'applicazione del filtro sia il suo \emph{cleanup}.
  \item[\emph{Pagination}] Legge la \emph{label} di range del \texttt{mat-paginator},
    clicca \emph{next page} e \emph{poll} sulla \emph{label} finché cambia, clicca
    \emph{prev page} e verifica che torni identica. Un secondo test cambia la
    \emph{page size} e verifica che il numero di righe si adegui.
  \item[\emph{Resize}] Misura la \texttt{width} dell'\emph{header} di una colonna,
    esegue un \texttt{mouse.down} sulla \texttt{.resize-holder}, una sequenza di
    \texttt{mouse.move} a passi intermedi (per simulare un \emph{drag} fluido che
    la \texttt{ResizeColumnsDirective} possa intercettare), un \texttt{mouse.up}, e
    asserisce che il delta assoluto sia $> 5$\,px.
  \item[\emph{Export Excel}] Registra una \texttt{Promise} su
    \texttt{page.waitForEvent('download')} \emph{prima} di cliccare il pulsante di
    esportazione (l'ordine è fondamentale: l'evento è sincrono rispetto al click e
    perderlo significa attendere all'infinito), poi clicca la voce ``Esporta in
    excel'' nel menu azioni rapide, e verifica che il \texttt{suggestedFilename}
    termini per \texttt{.xlsx}.
  \item[\emph{Refresh}] Registra una \texttt{Promise} su \texttt{page.waitForRequest}
    che si risolve quando parte una \texttt{POST} verso gli \emph{endpoint} della
    griglia, clicca il pulsante \emph{refresh}, e verifica che la richiesta sia
    partita entro $15$ secondi: il pattern \texttt{waitForRequest} osserva il
    traffico HTTP del \emph{browser} e dimostra che il \emph{refresh} non è solo un
    \emph{ri-render} cosmetico ma una vera nuova chiamata al \emph{backend}.
  \item[\emph{Multi-grid}] Naviga via URL diretto alla route
    \texttt{multi-grid/<id>/<guid>}, verifica che il \emph{host}
    \texttt{<sia-multi-grid>} sia visibile e che il numero di \emph{grid} figlie
    renderizzate sia $\geq 2$. Copre il \emph{wiring} del componente
    \texttt{sia-multi-grid}, area non esercitabile via Vitest senza render pesante.
  \item[\emph{Error 500}] Intercetta via \texttt{page.route} tutte le chiamate HTTP
    verso gli endpoint della griglia e le sostituisce con una risposta sintetica
    \texttt{500 Internal Server Error}, senza spegnere il \emph{backend}. Il test
    verifica che la UI non vada in \emph{crash} e che compaia una notifica di errore
    (\texttt{.dx-toast} o \texttt{.mat-mdc-snack-bar-container}). È il test che
    formalizza il contratto di robustezza: nessun errore HTTP deve portare
    l'interfaccia in uno stato non recuperabile.
\end{description}

Un aspetto tecnico rilevante è la gestione dell'autenticazione: eseguire il
\emph{login} via UI prima di ogni \emph{spec} sarebbe lento e ridondante. Il pattern
adottato prevede un \texttt{globalSetup} che esegue il \emph{login} una sola volta
all'avvio della suite, naviga a \texttt{/\#/auth/login}, compila i campi (letti da
\texttt{.env.e2e}, \emph{gitignored}) e salva lo \emph{storage state} (\emph{cookie}
+ \emph{localStorage} + \emph{sessionStorage}) su disco; tutti gli \emph{spec}
ereditano questo stato e partono già autenticati. Una sottigliezza scoperta sul
portale è che il \emph{token} di accesso viene memorizzato in
\texttt{sessionStorage}, che Playwright non persiste di default: la soluzione è
estrarlo manualmente al termine del \texttt{globalSetup} e re-iniettarlo tramite
\texttt{addInitScript} prima di ogni navigazione.

Le interazioni con la griglia sono incapsulate in tre \emph{page object}
(\texttt{LoginPage}, \texttt{GridPage}, \texttt{MultiGridPage}) sotto
\texttt{e2e/pages/}: ogni \emph{spec} consuma metodi semantici
(\texttt{gotoViaMenu}, \texttt{rowCount}, \texttt{clickHeaderToSort},
\texttt{setColumnFilter}, \texttt{resizeFirstColumn}, \dots) anziché \emph{locator}
crudi, isolando i selettori in un solo punto.

\section{Vantaggi ottenuti}

L'impatto più immediato dell'introduzione dei test è stata la scoperta di difetti
latenti che non sarebbero emersi prima di un incidente in produzione. Iterando
automaticamente su tutti i componenti del database, sono stati trovati due punti del
codice in cui certe combinazioni di dati causavano \emph{crash} con messaggi criptici:
un \texttt{NullReferenceException} in \texttt{CrudDataService.Get} su
\texttt{Filter} nullo, e una gestione difensiva mancante in \texttt{GridDataService}.
Entrambi sono stati corretti prima che il codice uscisse dall'ambiente di sviluppo.
Sul piano infrastrutturale, è stata resa visibile la situazione preesistente dei
$316$ \emph{synonym} obsoleti che mascheravano la non operatività dell'audit log.

Sul piano strategico, il beneficio più duraturo è la trasformazione del modello di
verifica da implicito (``funziona se non lo segnala nessuno'') a esplicito (``ogni
regressione che attraversa la suite viene catturata''). Modificare il framework con
fiducia, documentare il comportamento atteso in forma eseguibile, accorciare
l'\emph{onboarding} dei nuovi sviluppatori: sono conseguenze che si accumulano nel
tempo. Per garantire che questo investimento non si esaurisca con i moduli pilota, è
stata formalizzata una convenzione di scaffolding --- codificata nelle regole
degli agenti di sviluppo \texttt{sia-server-developer} e \texttt{sia-angular-developer}
--- per cui ogni nuovo modulo creato nel framework nasce con il proprio progetto di
test gemello.
